\documentclass[12pt,a4paper]{article}
\usepackage{mathwizards-ingles}
\title{%
  \vspace{-1cm}
  {\Huge\bfseries\textcolor{mineria}{Guía de Minería de Oraciones}}\\[0.3cm]
  {\Large Sentence Mining · La Columna Vertebral del Estudio Deliberado}\\[0.2cm]
  {\large\color{grissuave} Regla \textit{i+1} · Anatomía de la Tarjeta · Errores Comunes · Crowdsourcing}\\[0.5cm]
  \rule{\textwidth}{1.5pt}
}
\author{}
\date{}

\begin{document}
\maketitle
\thispagestyle{empty}

% ═══════════════════════════════════════════════════════════════════════════════
\section{¿Qué es la Minería de Oraciones?}
% ═══════════════════════════════════════════════════════════════════════════════

\begin{cajaprincipio}[Definición]
\textbf{Sentence Mining} (Minería de Oraciones) es el proceso de extraer oraciones del contenido que consumes (videos, podcasts, libros, páginas web) para crear tarjetas de memoria en Anki.

\medskip
A diferencia de memorizar listas de vocabulario aislado, la minería preserva el \textbf{contexto natural} de la palabra, lo que facilita su adquisición inconsciente y su uso posterior en producción.
\end{cajaprincipio}

% ═══════════════════════════════════════════════════════════════════════════════
\section{La Regla Fundamental: \textit{i+1} (Una Sola Incógnita)}
% ═══════════════════════════════════════════════════════════════════════════════

\begin{cajamineria}[El Principio de Aislamiento Cognitivo: La Regla ``1T'']
Cada oración minada debe contener \textbf{exactamente un elemento desconocido}. Todo lo demás en la oración debe ser comprensible.

\medskip
\begin{itemize}[leftmargin=1.2cm]
  \item \textbf{``i''} = tu nivel actual (lo que ya sabes)
  \item \textbf{``+1''} = un solo elemento nuevo (la incógnita)
\end{itemize}

\medskip
El cerebro puede deducir el significado del ``+1'' a partir del contexto conocido. Si hay \textbf{dos o más incógnitas}, el contexto se rompe y la adquisición falla.

\medskip
\begin{cajaejemplo}[Evaluación de una Oración]
\textbf{Tu nivel:} conoces ``I'', ``went'', ``to'', ``the'', ``yesterday''.

\medskip
$\checkmark$ \textit{I went to the \underline{supermarket} yesterday.}\\
{\small Solo una incógnita: ``supermarket''. El contexto permite deducirla.}

\medskip
$\times$ \textit{I \underline{strolled} to the \underline{quaint} \underline{boutique} yesterday.}\\
{\small Tres incógnitas. Imposible deducir las tres simultáneamente.}
\end{cajaejemplo}
\end{cajamineria}

% ═══════════════════════════════════════════════════════════════════════════════
\section{Anatomía de la Tarjeta Anki Perfecta}
% ═══════════════════════════════════════════════════════════════════════════════

\begin{cajaprincipio}[Estructura Estricta de la Tarjeta]
La tarjeta tiene \textbf{dos lados}. El diseño es rígido e innegociable para prevenir malas prácticas:

\medskip
\begin{cajatarjeta}[FRENTE (lo que ves primero)]
\begin{center}
\large
\textit{She has never \underline{\textbf{tried}} sushi, but she wants to.}
\end{center}
\begin{itemize}[nosep, leftmargin=0.8cm]
  \item La oración completa en inglés
  \item La palabra/expresión objetivo está \textbf{resaltada} (negrita, subrayado o color)
  \item \textbf{NO} hay traducción al español en este lado
  \item \textbf{NO} hay definición en este lado
\end{itemize}
\end{cajatarjeta}

\medskip
\begin{cajatarjeta}[REVERSO (lo que ves después de intentar recordar)]
\textbf{tried} /traɪd/ \quad 🔊 {\small [audio nativo]}

\medskip
\textbf{Definición (monolingüe):}\\
\textit{Past tense of ``try''. To taste or use something for the first time to see if you like it.}

\medskip
\textbf{Definición (bilingüe, solo Mes 1--2):}\\
\textit{Intentó, probó.}

\medskip
\textbf{Imagen (opcional):} [foto de alguien probando comida nueva]

\begin{itemize}[nosep, leftmargin=0.8cm]
  \item Definición clara (monolingüe a partir del Mes 2)
  \item Audio de un hablante nativo
  \item Imagen como anclaje visual (opcional pero recomendado)
\end{itemize}
\end{cajatarjeta}
\end{cajaprincipio}

% ═══════════════════════════════════════════════════════════════════════════════
\section{El Protocolo de Evaluación: ¿Cómo Calificar la Tarjeta?}
% ═══════════════════════════════════════════════════════════════════════════════

\begin{cajamineria}[¿Cuándo presionar cada botón?]
Al ver el frente de la tarjeta, \textbf{antes} de voltearla, pregúntate: ``¿Puedo entender esta oración \textbf{completamente}, incluyendo la palabra resaltada?''

\medskip
\begin{tabularx}{\textwidth}{@{} l l X @{}}
\toprule
\textbf{Botón} & \textbf{Criterio} & \textbf{Cuándo} \\
\midrule
\textcolor{bueno}{\textbf{Good}} & Entendí todo inmediatamente & La oración y la palabra objetivo fueron claras al instante \\
\textcolor{malo}{\textbf{Again}} & No recordé / no entendí & No pude recuperar el significado de la palabra objetivo \\
\textcolor{acento}{\textbf{Easy}} & Demasiado fácil & Si después de varias semanas la tarjeta es trivial \\
\textcolor{ejemplo}{\textbf{Hard}} & Difícil pero recordé & Tardé en recordar pero lo logré \\
\bottomrule
\end{tabularx}

\medskip
\begin{cajaejemplo}[Regla de Oro]
Si necesitas \textbf{traducir mentalmente al español} para entender, presiona \textbf{Again}. El objetivo es la comprensión \textbf{directa} en inglés, sin intermediarios.
\end{cajaejemplo}
\end{cajamineria}

% ═══════════════════════════════════════════════════════════════════════════════
\section{Errores Comunes (y Cómo Evitarlos)}
% ═══════════════════════════════════════════════════════════════════════════════

\begin{cajamalo}[Los 7 Pecados Capitales de la Minería]

\begin{enumerate}[leftmargin=1.2cm]
  \item \textbf{Poner la traducción en el FRENTE de la tarjeta.}\\
    {\small Esto convierte la tarjeta en un ejercicio de traducción, no de comprensión. El frente debe tener \textbf{solo inglés}.}

  \item \textbf{Minar oraciones con 2+ incógnitas.}\\
    {\small Si no entiendes al menos el 90\% de la oración, no la mines. Busca una más fácil.}

  \item \textbf{Agregar más de 10 tarjetas nuevas por día.}\\
    {\small Las tarjetas nuevas se acumulan como una bola de nieve. 10 tarjetas/día generan $\sim$70 repasos/día en 2 semanas.}

  \item \textbf{No suspender sanguijuelas (\textit{leeches}).}\\
    {\small Si fallas una tarjeta 5+ veces, suspéndela. Reformúlala con mejor contexto o elimínala.}

  \item \textbf{Minar palabras aisladas sin contexto.}\\
    {\small $\times$ Frente: ``\textit{apple}'' Reverso: ``manzana'' \\
    $\checkmark$ Frente: ``\textit{She always eats an \underline{apple} for breakfast.}''}

  \item \textbf{Minar vocabulario que nunca usarás.}\\
    {\small Prioriza palabras que aparecen en tu inmersión actual. Si una palabra aparece una sola vez en un libro, probablemente no vale la pena minarla.}

  \item \textbf{No revisar Anki diariamente.}\\
    {\small La repetición espaciada funciona \textbf{solo} si se respeta el calendario. Un día saltado arruina los intervalos.}
\end{enumerate}
\end{cajamalo}

% ═══════════════════════════════════════════════════════════════════════════════
\section{Flujo de Crowdsourcing Grupal}
% ═══════════════════════════════════════════════════════════════════════════════

\begin{cajabueno}[Minería Colaborativa: El Proceso Grupal]
En este curso, la minería tiene un componente \textbf{social}. El flujo es:

\medskip
\begin{enumerate}[leftmargin=1.2cm]
  \item \textbf{El estudiante mina una oración} durante su inmersión individual (video, lectura, podcast).
  
  \item \textbf{La publica en el canal grupal} (Discord/Slack) con el formato:
  \begin{center}
  \fbox{\parbox{0.8\textwidth}{
    \textbf{Oración:} \textit{She has never \underline{tried} sushi.}\\
    \textbf{Fuente:} Video de EnglishSponge, minuto 3:45\\
    \textbf{Mi definición:} To taste something for the first time.
  }}
  \end{center}
  
  \item \textbf{El profesor valida} que:
  \begin{itemize}[nosep, leftmargin=0.8cm]
    \item La oración tiene exactamente \textbf{una} incógnita (regla \textit{i+1})
    \item La definición es correcta
    \item La pronunciación/audio existe
  \end{itemize}
  
  \item \textbf{Otros estudiantes aprovechan} las oraciones validadas:
  \begin{itemize}[nosep, leftmargin=0.8cm]
    \item Si la incógnita también es nueva para ellos, la agregan a su mazo
    \item Si ya la conocen, la saltan
  \end{itemize}
  
  \item \textbf{El profesor detecta patrones:} Si múltiples estudiantes minan la misma palabra, la marca como \textbf{vocabulario de alta frecuencia} del grupo y la incluye en el mazo semanal compartido.
\end{enumerate}
\end{cajabueno}

% ═══════════════════════════════════════════════════════════════════════════════
\section{Transición: De Bilingüe a Monolingüe}
% ═══════════════════════════════════════════════════════════════════════════════

\begin{cajatransicion}[Cronograma de Transición a Definiciones Monolingües]
\begin{tabularx}{\textwidth}{@{} l l X @{}}
\toprule
\textbf{Periodo} & \textbf{Tipo de Definición} & \textbf{Detalle} \\
\midrule
Mes 1 (Sem 1--4) & 100\% bilingüe & Español $\rightarrow$ Inglés. Mazo compartido del profesor. \\
Mes 2 (Sem 5--6) & Bilingüe + monolingüe & Definición en inglés simple + traducción como backup. \\
Mes 2 (Sem 7--8) & Monolingüe prioritario & \textit{Learner's Dictionary}. Traducción solo si no se entiende la def. en inglés. \\
Mes 3 (Sem 9--12) & 100\% monolingüe & Solo definiciones en inglés. Sin traducción al español. \\
\bottomrule
\end{tabularx}

\medskip
\begin{cajaejemplo}[¿Qué diccionario monolingüe usar?]
Para aprendices, usar \textbf{Learner's Dictionaries}, que definen palabras con vocabulario controlado ($\sim$2000 palabras):
\begin{itemize}[nosep, leftmargin=0.8cm]
  \item Oxford Learner's Dictionary
  \item Cambridge Learner's Dictionary
  \item Longman Dictionary of Contemporary English
  \item Macmillan Dictionary
\end{itemize}
\textbf{NO} usar diccionarios nativos (Merriam-Webster, Oxford English Dictionary) hasta nivel B2+.
\end{cajaejemplo}
\end{cajatransicion}

% ═══════════════════════════════════════════════════════════════════════════════
\section{Parámetros de Anki (Recordatorio)}
% ═══════════════════════════════════════════════════════════════════════════════

\begin{cajaprincipio}[Configuración Recomendada]
\begin{tabularx}{\textwidth}{@{} l X @{}}
\toprule
\textbf{Parámetro} & \textbf{Valor} \\
\midrule
Algoritmo & \textbf{FSRS} (no SM-2) \\
Tarjetas nuevas/día & \textbf{5--10} (nunca más de 10) \\
Repasos máximos/día & \textbf{Sin límite} (9999) \\
Retención deseada & \textbf{0.90} (90\%) \\
Umbral de sanguijuela & \textbf{5 fallos} $\rightarrow$ suspender automáticamente \\
\bottomrule
\end{tabularx}

\medskip
\textbf{Tiempo diario estimado en Anki:} 10--20 minutos (repasos) + 5--10 minutos (agregar tarjetas nuevas).
\end{cajaprincipio}

\vfill
\begin{center}
\rule{0.5\textwidth}{0.5pt}\\[0.3cm]
{\small\color{grissuave} Guía de Minería de Oraciones · Metodología Refold · Ecosistema Anki + FSRS}
\end{center}

\end{document}
